%----------------------------------------------------------
% 제5장. AI 정책 및 리스크 관리 체계
%----------------------------------------------------------

\chapter{AI 정책 및 리스크 관리 체계}
\label{ch:policy_risk}

AI 거버넌스의 실효성은 명확하고 체계적인 정책과 이를 뒷받침하는 리스크 관리 체계에서 시작된다. 정책은 조직의 AI 활용 원칙과 기준을 공식화\cite{schneider2022governance}하고 구성원의 행동 지침을 제공하는 '헌법'과 같으며, 리스크 관리는 예측 불가능한 AI의 특성으로부터 조직을 보호하는 '방어선' 역할을 한다.

본 장에서는 AI 정책 체계의 설계, 핵심 정책 문서 개발, AI 영향평가\cite{raji2020closing}(AI Impact Assessment) 체계, NIST AI RMF\cite{iso23894}를 준수한 리스크 식별 및 평가 방법론, 정책 준수\cite{mokander2022conformity} 모니터링\cite{klaise2021monitoring}, Policy-as-Code\cite{hinrichs2019opa} 자동화, 그리고 인시던트 관리\cite{brundage2020trustworthy} 및 위기 대응 체계를 통합적으로 다룬다. 이를 통해 조직이 AI 거버넌스를 문서 수준에서 운영 수준으로 전환할 수 있는 실질적 가이드라인을 제공한다.

%----------------------------------------------------------
\section{정책 아키텍처 및 계층 구조}
\label{sec:5.1}

\subsection{정책 체계의 3계층 구조}
\label{sec:5.1.1}

효과적인 AI 정책 체계는 계층적 구조를 가진다. 이는 ISO/IEC 42001(AI 경영시스템)의 구조와도 일맥상통한다 \cite{iso42001}. 최상위 정책은 조직의 AI 활용에 대한 원칙과 방향성을 선언하고, 중간 계층의 표준과 지침은 구체적인 요구사항을 정의하며, 최하위 계층의 절차와 매뉴얼은 실무자가 바로 적용할 수 있는 상세 지침을 제공한다.

\begin{figure}[htbp]
\centering
\begin{tikzpicture}[
 level/.style={rectangle, draw, rounded corners, minimum width=10cm, minimum height=1.2cm, align=center, font=\small},
]
\node[level, fill=blue!10, draw=blue!60, thick] (l1) at (0, 4) {\textbf{1계층: 정책(Policy)}\\원칙, 가치, 방향성 선언};
\node[level, fill=green!10, draw=green!60, thick] (l2) at (0, 2) {\textbf{2계층: 표준/지침(Standard/Guideline)}\\구체적 요구사항, 기준, 절차};
\node[level, fill=orange!10, draw=orange!60, thick] (l3) at (0, 0) {\textbf{3계층: 절차/매뉴얼(Procedure/Manual)}\\상세 실행 방법, 템플릿, 도구};
\draw[->, thick, >=stealth] (l1) -- (l2);
\draw[->, thick, >=stealth] (l2) -- (l3);
\end{tikzpicture}
\caption{AI 정책 체계 3계층 구조}
\label{fig:policy_hierarchy}
\end{figure}

\begin{table}[htbp]
\centering
\caption{AI 정책 계층별 특성}
\label{tab:policy_layers}
\begin{tabularx}{\textwidth}{p{1.5cm}p{2.5cm}p{2.5cm}X}
\toprule
\textbf{계층} & \textbf{문서 유형} & \textbf{승인 권한} & \textbf{내용 범위} \\
\midrule
1계층 & 정책(Policy) & 이사회, CEO & 조직의 AI 활용 원칙, 전략적 방향 \\
2계층 & 표준, 지침 & CDO, CTO & 준수해야 할 구체적 요구사항, 기준 \\
3계층 & 절차, 매뉴얼 & 부서장 & 상세 실행 방법, 템플릿, 도구 사용법 \\
\bottomrule
\end{tabularx}
\end{table}

\subsection{1계층: 정책 문서 상세 템플릿}
\label{sec:5.1.2}

1계층 정책은 조직의 AI 활용에 대한 최상위 방향성을 선언하는 문서로서, 이사회 또는 최고경영자의 승인을 필요로 한다. 일반적으로 세 가지 핵심 정책 문서로 구성된다.

\paragraph{AI 윤리 정책(AI Ethics Policy)}
AI 윤리 정책은 조직이 AI를 개발하고 활용함에 있어 준수해야 할 윤리적 원칙과 가치를 공식 선언하는 문서이다. 이 정책은 공정성\cite{mehrabi2021survey}, 투명성\cite{arrieta2020explainable}, 책임성, 안전성, 프라이버시 보호 등 핵심 원칙을 포함해야 한다.

\begin{table}[htbp]
\centering
\caption{AI 윤리 정책 템플릿 구성}
\label{tab:ethics_policy_template}
\begin{tabularx}{\textwidth}{p{3cm}X}
\toprule
\textbf{항목} & \textbf{내용} \\
\midrule
정책 목적 & AI 활용에 대한 윤리적 원칙과 기준 수립 \\
적용 범위 & 조직 내 모든 AI 시스템 및 관련 활동 \\
핵심 원칙 & 공정성, 투명성, 책임성, 안전성, 프라이버시 \\
역할과 책임 & AI 윤리위원회, 개발팀, 운영팀 등 역할 정의 \\
위반 시 조치 & 위반 유형별 제재 수준 및 절차 \\
검토 주기 & 연 1회 정기 검토, 중대 변경 시 수시 검토 \\
관련 문서 & 하위 표준, 지침, 절차 문서 목록 \\
\bottomrule
\end{tabularx}
\end{table}

\paragraph{AI 활용 정책(AI Use Policy)}
AI 활용 정책은 조직 내에서 AI 기술을 어떤 목적으로, 어떤 범위에서, 어떤 조건하에 활용할 수 있는지를 규정한다. 허용 가능한 활용 사례와 금지 사례를 명확히 구분하고, 생성형 AI(Generative AI) 활용에 대한 별도 조항\cite{bender2021stochastic}을 포함하는 것이 바람직하다.

\begin{itemize}
 \item \textbf{허용 활용 사례}: 업무 효율화, 데이터 분석, 고객 서비스 자동화 등
 \item \textbf{조건부 허용}: 고위험 의사결정 지원(인적 검토 필수 조건)
 \item \textbf{금지 사례}: 무기 개발, 불법 감시, 차별적 의사결정 자동화
 \item \textbf{생성형 AI 조항}: 기밀 정보 입력 금지, 산출물 검증 의무, 저작권 준수
\end{itemize}

\paragraph{AI 보안 정책(AI Security Policy)}
AI 보안 정책은 AI 시스템의 전체 생명주기(Lifecycle)에 걸친 보안 요구사항을 정의한다. 모델 보안, 데이터 보안, 인프라 보안, 적대적 공격 방어를 포괄하며, 기존 정보보호 정책과의 연계를 명시\cite{biggio2018wild}해야 한다.

\begin{table}[htbp]
\centering
\caption{AI 보안 정책 핵심 영역}
\label{tab:security_policy_areas}
\begin{tabularx}{\textwidth}{p{2.5cm}p{3cm}X}
\toprule
\textbf{영역} & \textbf{보호 대상} & \textbf{주요 통제 항목} \\
\midrule
모델 보안 & 학습 모델, 가중치 & 모델 접근 통제, 모델 추출 방지, 버전 관리 \\
데이터 보안 & 학습/추론 데이터 & 암호화, 접근 권한 관리, 데이터 유출 방지 \\
인프라 보안 & 서버, GPU 클러스터 & 네트워크 분리, 취약점 관리, 로그 감시 \\
적대적 방어 & AI 시스템 전체 & 입력 검증, 적대적 학습, 이상 탐지 \\
공급망 보안 & 외부 모델, 라이브러리 & 서드파티 검증, 종속성 관리, SBOM \\
\bottomrule
\end{tabularx}
\end{table}

\subsection{2계층: 표준 및 지침}
\label{sec:5.1.3}

2계층 문서는 1계층 정책을 구체화하여 실무에서 준수해야 할 요구사항과 기준을 제시한다. 주요 표준 및 지침은 다음과 같다.

\paragraph{모델 개발 표준(Model Development Standard)}
모델 개발 표준은 AI 모델의 설계, 학습, 검증, 배포에 이르는 전 과정에서 준수해야 할 기술적 기준을 정의한다.

\begin{enumerate}
 \item \textbf{설계 단계}: 문제 정의서 작성, 적합 알고리즘 선정 기준, 편향 리스크 사전 평가
 \item \textbf{데이터 준비}: 데이터 품질 기준(완전성, 정확성, 적시성), 편향 탐지 절차
 \item \textbf{학습 단계}: 하이퍼파라미터 관리, 실험 추적, 재현성 보장
 \item \textbf{검증 단계}: 성능 기준치(Threshold) 충족, 공정성 메트릭 검증, 강건성 테스트
 \item \textbf{배포 단계}: 배포 승인 게이트, 카나리 배포(Canary Deployment), 롤백 계획
\end{enumerate}

\paragraph{데이터 활용 지침(Data Use Guideline)}
데이터 활용 지침은 AI 시스템에 사용되는 데이터의 수집, 저장, 처리, 공유에 관한 원칙과 절차를 규정한다. 개인정보보호법 및 데이터 3법과의 정합성을 확보해야 한다.

\begin{table}[htbp]
\centering
\caption{데이터 활용 지침 핵심 항목}
\label{tab:data_guideline}
\begin{tabularx}{\textwidth}{p{2.5cm}X}
\toprule
\textbf{항목} & \textbf{요구사항} \\
\midrule
데이터 수집 & 합법적 근거 확보, 동의 관리, 목적 제한 원칙 준수 \\
데이터 품질 & 정확성, 완전성, 일관성 검증 절차 수립 \\
데이터 보관 & 보관 기간 설정, 기간 경과 시 삭제 또는 비식별화 \\
데이터 공유 & 내부 공유 승인 절차, 외부 공유 시 계약 및 기술적 보호조치 \\
민감 데이터 & 특별 카테고리 데이터(인종, 건강 등) 처리 제한 \\
합성 데이터 & 합성 데이터 생성 및 활용 기준, 원본 데이터 유출 방지 \\
\bottomrule
\end{tabularx}
\end{table}

\paragraph{설명가능성 지침(Explainability Guideline)}
설명가능성(Explainability) 지침은 AI 시스템의 의사결정 과정을 이해관계자가 이해할 수 있도록 설명하는 방법과 수준을 규정한다. EU AI Act와 한국 AI 기본법\cite{cset2025korea_ai_act} 모두 고위험 AI 시스템에 대한 설명가능성을 요구하고 있다.

\begin{itemize}
 \item \textbf{설명 수준}: 기술적 설명(개발자용), 운영적 설명(운영자용), 사용자 설명(최종 사용자용)
 \item \textbf{설명 방법}: 전역 설명(Global Explanation), 지역 설명(Local Explanation)
 \item \textbf{설명 기법}: SHAP, LIME, 주의(Attention) 시각화, 반사실적 설명(Counterfactual)
 \item \textbf{문서화}: 모델 카드(Model Card) 작성, 설명서 템플릿 준수
 \item \textbf{검증}: 설명의 충실도(Fidelity) 측정, 사용자 이해도 테스트
\end{itemize}

\subsection{3계층: 절차 및 매뉴얼}
\label{sec:5.1.4}

3계층 문서는 실무자가 구체적인 업무 상황에서 직접 참조하여 실행할 수 있는 상세 절차와 매뉴얼이다. 부서장 수준에서 승인되며, 필요에 따라 수시로 갱신된다.

\begin{table}[htbp]
\centering
\caption{3계층 절차 매뉴얼 예시}
\label{tab:procedure_examples}
\begin{tabularx}{\textwidth}{p{3.5cm}p{2cm}X}
\toprule
\textbf{매뉴얼명} & \textbf{대상} & \textbf{주요 내용} \\
\midrule
모델 배포 승인 절차서 & 개발팀 & 배포 전 체크리스트, 승인 워크플로, 승인자 목록 \\
데이터 비식별화 매뉴얼 & 데이터팀 & 비식별화 기법별 적용 방법, 재식별 리스크 평가 \\
인시던트 대응 매뉴얼 & 운영팀 & 인시던트 분류, 에스컬레이션 경로, 보고 양식 \\
공정성 검증 절차서 & QA팀 & 공정성 메트릭 산출, 임계값 판단, 개선 조치 \\
생성형 AI 활용 매뉴얼 & 전 직원 & 허용/금지 행위, 프롬프트 작성법, 산출물 검증 \\
외부 API 연동 절차서 & 개발팀 & 외부 AI 서비스 평가, 계약 요건, 보안 요구사항 \\
\bottomrule
\end{tabularx}
\end{table}

\subsection{정책 수립 프로세스}
\label{sec:5.1.5}

AI 정책은 일회성 문서가 아니라 지속적으로 발전해야 하는 살아있는 체계이다. 효과적인 정책 수립을 위해서는 체계적인 프로세스가 필요하다.

\paragraph{이해관계자 참여(Stakeholder Engagement)}
정책의 실효성은 다양한 이해관계자의 참여에 의해 좌우된다. 정책 수립 시 반드시 다음 이해관계자의 의견을 수렴해야 한다.

\begin{itemize}
 \item \textbf{경영진}: 전략적 방향, 리스크 허용 수준(Risk Appetite) 결정
 \item \textbf{법무/컴플라이언스}: 법적 요구사항 반영, 규제 동향 분석
 \item \textbf{AI 개발팀}: 기술적 실현가능성 검토, 현장 적용성 확인
 \item \textbf{사업부}: 비즈니스 영향 평가, 활용 사례 검토
 \item \textbf{데이터 보호 책임자(DPO)}: 개인정보 영향 평가, 프라이버시 요구사항
 \item \textbf{외부 전문가}: 독립적 시각, 업계 모범 사례 자문
\end{itemize}

\paragraph{검토 주기 및 변경 관리}
정책은 정기적으로 검토되어야 하며, 환경 변화에 따른 수시 검토 메커니즘도 갖추어야 한다.

\begin{table}[htbp]
\centering
\caption{정책 검토 주기 및 변경 관리 기준}
\label{tab:policy_review_cycle}
\begin{tabularx}{\textwidth}{p{2.5cm}p{2cm}X}
\toprule
\textbf{검토 유형} & \textbf{주기} & \textbf{트리거 조건} \\
\midrule
정기 검토 & 연 1회 & 회계연도 종료 후 1분기 내 \\
수시 검토 & 비정기 & 신규 법규 시행, 중대 인시던트 발생, 신기술 도입 \\
긴급 검토 & 즉시 & 법적 제재, 심각한 보안 사고, 규제 기관 권고 \\
\bottomrule
\end{tabularx}
\end{table}

변경 관리 프로세스는 다음 단계를 따른다: (1) 변경 요청서 제출, (2) 영향 분석 수행, (3) 이해관계자 검토 및 의견 수렴, (4) 승인 권한자 결재, (5) 변경 내용 공지 및 교육, (6) 시행 및 모니터링.

\paragraph{정책 거버넌스 조직과의 연계}
정책 체계는 제4장에서 다룬 거버넌스 조직 체계와 긴밀히 연계되어야 한다(\ref{ch:governance_organization}장 참조). AI 윤리위원회는 1계층 정책의 최종 승인과 감독을, AI 거버넌스팀(COE)은 2계층 표준 및 지침의 개발과 유지를, 각 사업부 AI 담당자는 3계층 절차의 현장 적용과 피드백을 담당한다. 이러한 역할 분담이 명확할 때 정책은 선언에 그치지 않고 실제로 조직 내에서 작동하게 된다.

%----------------------------------------------------------
\section{AI 영향평가 체계}
\label{sec:5.2}

AI 영향평가(AI Impact Assessment)는 AI 시스템이 개인, 집단, 사회에 미칠 수 있는 영향을 사전에 식별하고 평가하여 필요한 완화 조치를 수립하는 체계적 절차이다. EU AI Act과 한국 AI 기본법 모두 고위험 AI 시스템에 대한 영향평가를 의무화하고 있어, 조직 차원의 체계 구축이 필수적이다 \cite{eu2024aiact}.

\subsection{EU AI Act 기반 고위험 AI 영향평가}
\label{sec:5.2.1}

EU AI Act는 기본권 영향평가(Fundamental Rights Impact Assessment, FRIA)를 고위험 AI 시스템 운용 기관에 의무화하고 있다. 주요 평가 항목은 다음과 같다.

\begin{enumerate}
 \item \textbf{시스템 설명}: AI 시스템의 목적, 기능, 범위 기술
 \item \textbf{기본권 영향}: 차별 금지, 프라이버시, 표현의 자유, 사법 접근성 등에 대한 영향
 \item \textbf{취약 계층 영향}: 아동, 장애인, 고령자 등 취약 집단에 대한 특별 고려
 \item \textbf{인적 감독 조치}: 자동화된 의사결정에 대한 인간 개입 장치
 \item \textbf{완화 조치}: 식별된 리스크에 대한 기술적, 조직적 대응 방안
 \item \textbf{모니터링 계획}: 배포 후 지속적 영향 추적 방안
\end{enumerate}

\subsection{한국 AI 기본법 영향평가 요구사항}
\label{sec:5.2.2}

한국 AI 기본법(2025년 시행)은 고위험 AI 시스템에 대한 사전 영향평가를 요구한다. 주요 요구사항은 다음과 같다.

\begin{table}[htbp]
\centering
\caption{한국 AI 기본법 영향평가 요구사항}
\label{tab:korea_ai_act_ia}
\begin{tabularx}{\textwidth}{p{3cm}X}
\toprule
\textbf{요구사항} & \textbf{내용} \\
\midrule
평가 대상 & 고위험 AI: 생명, 안전, 기본권에 중대한 영향을 미치는 AI \\
평가 시점 & AI 시스템 도입 전, 중대 변경 시 \\
평가 항목 & 안전성, 신뢰성, 공정성, 투명성, 책임성 \\
평가 주체 & AI 시스템 운용 기관(필요시 외부 전문기관 위탁) \\
결과 보고 & 과학기술정보통신부에 영향평가 결과 보고 \\
공개 의무 & 이용자에게 고위험 AI 사용 사실 및 평가 결과 공개 \\
\bottomrule
\end{tabularx}
\end{table}

\subsection{영향평가 단계별 절차}
\label{sec:5.2.3}

AI 영향평가는 4단계 절차를 통해 체계적으로 수행된다.

\paragraph{1단계: 사전 스크리닝(Pre-screening)}
사전 스크리닝은 해당 AI 시스템이 상세 영향평가 대상인지를 판별하는 초기 단계이다. 리스크 수준에 따라 최소 리스크(Minimal Risk), 제한 리스크(Limited Risk), 고위험(High Risk), 금지(Unacceptable Risk)로 분류한다.

\paragraph{2단계: 상세 평가(Detailed Assessment)}
사전 스크리닝에서 고위험으로 분류된 AI 시스템은 상세 영향평가를 수행한다. 기술적 평가(성능, 공정성, 강건성), 윤리적 평가(편향, 차별, 프라이버시), 사회적 평가(고용, 불평등, 사회적 영향)를 포괄적으로 수행한다.

\paragraph{3단계: 완화 조치(Mitigation Measures)}
상세 평가에서 식별된 리스크에 대해 기술적 조치(편향 제거, 설명가능성 강화), 조직적 조치(인적 감독 강화, 교육 시행), 절차적 조치(이의 제기 절차, 구제 메커니즘)를 수립한다.

\paragraph{4단계: 모니터링(Monitoring)}
완화 조치 시행 후 지속적으로 효과를 모니터링하고, 새로운 리스크가 발견되면 평가를 재수행한다.

\begin{table}[htbp]
\centering
\caption{영향평가 단계별 핵심 활동}
\label{tab:ia_stages}
\begin{tabularx}{\textwidth}{p{1cm}p{2.5cm}p{2.5cm}X}
\toprule
\textbf{단계} & \textbf{활동} & \textbf{산출물} & \textbf{주요 참여자} \\
\midrule
1 & 사전 스크리닝 & 리스크 분류 결과 & AI 거버넌스팀, 사업부 \\
2 & 상세 평가 & 영향평가 보고서 & 개발팀, 법무팀, 외부 전문가 \\
3 & 완화 조치 & 완화 계획서 & 개발팀, 운영팀, 거버넌스팀 \\
4 & 모니터링 & 모니터링 보고서 & 운영팀, 거버넌스팀 \\
\bottomrule
\end{tabularx}
\end{table}

\subsection{영향평가 자동화 도구}
\label{sec:5.2.4}

영향평가의 일관성과 효율성을 높이기 위해 자동화 도구를 활용할 수 있다. 다음은 사전 스크리닝부터 상세 평가까지 자동화하는 파이썬(Python) 구현 예시이다.

\begin{lstlisting}[language=Python, caption={AI 영향평가 자동화 도구}, label={lst:impact_assessment}]
from dataclasses import dataclass, field
from enum import Enum
from typing import Dict, List, Optional
from datetime import datetime


class RiskLevel(Enum):
    """EU AI Act 기반 위험 수준 분류"""
    UNACCEPTABLE = 4  # 금지
    HIGH = 3          # 고위험
    LIMITED = 2       # 제한 위험
    MINIMAL = 1       # 최소 위험


@dataclass
class AISystemProfile:
    """AI 시스템 프로파일"""
    system_id: str
    name: str
    purpose: str
    domain: str
    data_types: List[str]
    affected_population: List[str]
    automation_level: str  # full, partial, advisory
    decision_impact: str   # critical, significant, moderate, low


@dataclass
class ImpactAssessmentResult:
    """영향평가 결과"""
    system_id: str
    risk_level: RiskLevel
    scores: Dict[str, float]
    findings: List[str]
    mitigations: List[str]
    assessed_at: str = field(
        default_factory=lambda: datetime.now().isoformat()
    )


class AIImpactAssessor:
    """AI 영향평가 자동화 엔진"""

    # 고위험 도메인 (EU AI Act Annex III 기반)
    HIGH_RISK_DOMAINS = [
        "employment", "credit_scoring", "education",
        "law_enforcement", "immigration", "healthcare",
        "critical_infrastructure", "judiciary"
    ]

    # 민감 데이터 유형
    SENSITIVE_DATA_TYPES = [
        "biometric", "health", "financial",
        "criminal", "political", "ethnic", "religious"
    ]

    def pre_screening(
        self, profile: AISystemProfile
    ) -> RiskLevel:
        """1단계: 사전 스크리닝"""
        score = 0

        # 도메인 기반 평가
        if profile.domain in self.HIGH_RISK_DOMAINS:
            score += 3

        # 민감 데이터 사용 여부
        sensitive_count = sum(
            1 for dt in profile.data_types
            if dt in self.SENSITIVE_DATA_TYPES
        )
        score += min(sensitive_count, 3)

        # 자동화 수준
        automation_scores = {
            "full": 3, "partial": 2, "advisory": 1
        }
        score += automation_scores.get(
            profile.automation_level, 0
        )

        # 의사결정 영향도
        impact_scores = {
            "critical": 3, "significant": 2,
            "moderate": 1, "low": 0
        }
        score += impact_scores.get(
            profile.decision_impact, 0
        )

        # 위험 수준 분류
        if score >= 10:
            return RiskLevel.UNACCEPTABLE
        elif score >= 7:
            return RiskLevel.HIGH
        elif score >= 4:
            return RiskLevel.LIMITED
        else:
            return RiskLevel.MINIMAL

    def detailed_assessment(
        self,
        profile: AISystemProfile,
        fairness_score: float,
        transparency_score: float,
        safety_score: float,
        privacy_score: float,
        accountability_score: float
    ) -> ImpactAssessmentResult:
        """2단계: 상세 영향평가"""
        risk_level = self.pre_screening(profile)

        scores = {
            "fairness": fairness_score,
            "transparency": transparency_score,
            "safety": safety_score,
            "privacy": privacy_score,
            "accountability": accountability_score,
        }
        avg = sum(scores.values()) / len(scores)
        scores["overall"] = round(avg, 2)

        findings = []
        mitigations = []

        if fairness_score < 0.7:
            findings.append(
                "공정성 점수 미달: 편향 위험 존재"
            )
            mitigations.append(
                "편향 탐지 및 완화 기법 적용 필요"
            )
        if transparency_score < 0.6:
            findings.append(
                "투명성 부족: 설명가능성 개선 필요"
            )
            mitigations.append(
                "SHAP/LIME 기반 설명 모듈 추가"
            )
        if privacy_score < 0.7:
            findings.append(
                "프라이버시 위험: 데이터 보호 강화 필요"
            )
            mitigations.append(
                "차분 프라이버시 또는 연합학습 적용"
            )

        return ImpactAssessmentResult(
            system_id=profile.system_id,
            risk_level=risk_level,
            scores=scores,
            findings=findings,
            mitigations=mitigations,
        )


# 사용 예시
if __name__ == "__main__":
    assessor = AIImpactAssessor()

    profile = AISystemProfile(
        system_id="AI-2026-001",
        name="AI 채용 심사 시스템",
        purpose="이력서 자동 스크리닝 및 적합도 평가",
        domain="employment",
        data_types=["biometric", "ethnic"],
        affected_population=["job_applicants"],
        automation_level="partial",
        decision_impact="significant",
    )

    result = assessor.detailed_assessment(
        profile,
        fairness_score=0.65,
        transparency_score=0.55,
        safety_score=0.85,
        privacy_score=0.60,
        accountability_score=0.80,
    )

    print(f"위험 수준: {result.risk_level.name}")
    print(f"종합 점수: {result.scores['overall']}")
    for f in result.findings:
        print(f"  발견사항: {f}")
    for m in result.mitigations:
        print(f"  완화 조치: {m}")
\end{lstlisting}

\begin{practicaltool}{AI 영향평가 체크리스트}
다음 체크리스트는 AI 시스템 도입 전 영향평가를 수행할 때 활용할 수 있다.

\begin{enumerate}
 \item \textbf{시스템 식별}
 \begin{itemize}
  \item[$\square$] AI 시스템의 목적과 기능을 명확히 정의하였는가?
  \item[$\square$] 영향을 받는 이해관계자 집단을 식별하였는가?
  \item[$\square$] 자동화 수준과 인적 감독 메커니즘을 정의하였는가?
 \end{itemize}

 \item \textbf{리스크 분류}
 \begin{itemize}
  \item[$\square$] EU AI Act 또는 한국 AI 기본법 기준 리스크 수준을 분류하였는가?
  \item[$\square$] 고위험으로 분류된 경우 상세 평가 절차를 진행하였는가?
 \end{itemize}

 \item \textbf{공정성 평가}
 \begin{itemize}
  \item[$\square$] 보호 속성(성별, 연령, 인종 등)별 성능 차이를 측정하였는가?
  \item[$\square$] 통계적 패리티(Statistical Parity), 균등 기회(Equal Opportunity) 등 공정성 메트릭을 산출하였는가?
 \end{itemize}

 \item \textbf{투명성 및 설명가능성}
 \begin{itemize}
  \item[$\square$] 모델의 의사결정 과정을 이해관계자에게 설명할 수 있는가?
  \item[$\square$] 모델 카드(Model Card) 또는 데이터시트(Datasheet)를 작성하였는가?
 \end{itemize}

 \item \textbf{프라이버시 및 보안}
 \begin{itemize}
  \item[$\square$] 개인정보 영향평가(PIA)를 수행하였는가?
  \item[$\square$] 데이터 최소화 원칙을 준수하고 있는가?
 \end{itemize}

 \item \textbf{완화 조치 및 모니터링}
 \begin{itemize}
  \item[$\square$] 식별된 모든 리스크에 대한 완화 조치를 수립하였는가?
  \item[$\square$] 배포 후 모니터링 계획을 수립하였는가?
  \item[$\square$] 이의 제기 및 구제 메커니즘을 마련하였는가?
 \end{itemize}
\end{enumerate}
\end{practicaltool}

%----------------------------------------------------------
\section{AI 리스크 분류 및 평가}
\label{sec:5.3}

\subsection{5대 리스크 카테고리}
\label{sec:5.3.1}

AI 시스템의 리스크는 ``데이터가 곧 로직이 되는'' 특성 때문에 전통적인 IT 시스템과 근본적 차이를 보인다 \cite{nist2023airmf}. 조직은 AI 고유의 리스크를 포괄적으로 관리하기 위해 다섯 가지 핵심 카테고리를 정의하고, 각 카테고리별로 구체적인 리스크 항목을 식별해야 한다.

\begin{table}[htbp]
\centering
\caption{AI 5대 리스크 카테고리 상세}
\label{tab:risk_categories_detail}
\begin{tabularx}{\textwidth}{p{2cm}p{3.5cm}X}
\toprule
\textbf{카테고리} & \textbf{주요 리스크} & \textbf{구체적 사례} \\
\midrule
기술 리스크 & 성능 저하, 적대적 공격 & 모델 정확도 급락, 프롬프트 인젝션, 데이터 오염 공격, 모델 역공학 \\
\midrule
운영 리스크 & 데이터 드리프트, 파이프라인 장애 & 입력 분포 변화로 인한 예측 오류, MLOps 장애, 모니터링 사각지대 \\
\midrule
컴플라이언스 & 법규 위반, 인증 미비 & AI 기본법 미준수, GDPR 위반, 영향평가 미실시, 설명의무 불이행 \\
\midrule
전략 리스크 & 평판 하락, 가치 훼손 & AI 오작동 언론 보도, 고객 신뢰 상실, 비즈니스 가치 미달성 \\
\midrule
윤리 리스크 & 편향, 차별, 프라이버시 & 채용 AI 성별 편향, 대출 심사 인종 차별, 과도한 개인정보 수집 \\
\bottomrule
\end{tabularx}
\end{table}

\subsection{NIST AI RMF 4대 기능}
\label{sec:5.3.2}

NIST AI 리스크 관리 프레임워크(AI Risk Management Framework, AI RMF 1.0)는 AI 리스크를 체계적으로 관리하기 위한 4대 핵심 기능을 정의한다 \cite{nist2023airmf}. 각 기능은 독립적으로 작동하면서도 상호 보완적 관계를 형성한다.

\paragraph{거버넌스(Govern)}
거버넌스 기능은 조직의 AI 리스크 관리 문화, 정책, 프로세스를 수립하는 기반 기능이다. 다른 세 기능(Map, Measure, Manage)의 효과적 운영을 위한 전제 조건이며, 조직 전반에 걸쳐 횡단적으로 작동한다.

\begin{itemize}
 \item 리스크 관리 정책 및 절차 수립
 \item 역할과 책임 정의(RACI 매트릭스)
 \item 리스크 허용 수준(Risk Appetite) 결정
 \item 조직 문화 및 인식 제고
 \item 다양성 및 포용성 확보
 \item 서드파티 리스크 관리 체계
\end{itemize}

\paragraph{매핑(Map)}
매핑 기능은 AI 시스템의 맥락(Context)을 파악하고, 잠재적 리스크를 식별하는 단계이다. 시스템이 사용되는 환경, 영향을 받는 이해관계자, 예상되는 리스크 시나리오를 종합적으로 파악한다.

\begin{itemize}
 \item AI 시스템 목적과 사용 맥락 정의
 \item 이해관계자 식별 및 영향 분석
 \item 리스크 시나리오 도출
 \item 편향 원인 사전 분석
 \item 데이터 및 모델 한계 문서화
\end{itemize}

\paragraph{측정(Measure)}
측정 기능은 식별된 리스크를 정량적, 정성적으로 평가하고 추적하는 단계이다. 메트릭 정의, 테스트 수행, 결과 분석을 포함한다.

\begin{itemize}
 \item 리스크 평가 메트릭 정의
 \item 공정성, 정확성, 강건성 측정
 \item 설명가능성 평가
 \item 리스크 추적 및 트렌드 분석
 \item 독립적 평가 수행
\end{itemize}

\paragraph{관리(Manage)}
관리 기능은 측정 결과에 기반하여 리스크를 처리하고, 잔여 리스크를 지속적으로 모니터링하는 단계이다.

\begin{itemize}
 \item 리스크 대응 조치 결정 및 실행
 \item 리스크 우선순위화 및 자원 배분
 \item 인시던트 대응 및 복구
 \item 잔여 리스크 모니터링
 \item 지속적 개선
\end{itemize}

\begin{table}[htbp]
\centering
\caption{NIST AI RMF 4대 기능 요약}
\label{tab:nist_functions}
\begin{tabularx}{\textwidth}{p{2cm}p{2.5cm}p{3cm}X}
\toprule
\textbf{기능} & \textbf{핵심 질문} & \textbf{주요 활동} & \textbf{산출물} \\
\midrule
거버넌스 & 리스크 관리를 위한 조직 기반이 갖추어져 있는가? & 정책 수립, 역할 정의, 문화 조성 & 리스크 관리 정책, RACI 매트릭스 \\
\midrule
매핑 & 어떤 리스크가 존재하는가? & 맥락 분석, 이해관계자 식별, 리스크 시나리오 도출 & 리스크 목록, 영향 분석 보고서 \\
\midrule
측정 & 리스크의 크기는 얼마인가? & 메트릭 산출, 테스트, 분석 & 리스크 점수, 평가 보고서 \\
\midrule
관리 & 리스크를 어떻게 처리할 것인가? & 대응 실행, 모니터링, 개선 & 완화 계획서, 모니터링 보고서 \\
\bottomrule
\end{tabularx}
\end{table}

\subsection{리스크 매트릭스: 발생 가능성 $\times$ 영향도}
\label{sec:5.3.3}

리스크의 우선순위를 결정하기 위해 발생 가능성(Likelihood)과 영향도(Impact)를 교차 평가하는 리스크 매트릭스를 활용한다. 각 축은 5점 척도로 평가하며, 두 점수의 곱으로 리스크 등급을 산정한다.

\begin{table}[htbp]
\centering
\caption{리스크 매트릭스 (발생 가능성 $\times$ 영향도)}
\label{tab:risk_matrix}
\begin{tabularx}{\textwidth}{p{2.5cm}XXXXX}
\toprule
 & \textbf{영향도 1} & \textbf{영향도 2} & \textbf{영향도 3} & \textbf{영향도 4} & \textbf{영향도 5} \\
 & (미미) & (경미) & (보통) & (심각) & (치명) \\
\midrule
\textbf{가능성 5} (확실) & 5 & 10 & 15 & 20 & 25 \\
\textbf{가능성 4} (높음) & 4 & 8 & 12 & 16 & 20 \\
\textbf{가능성 3} (보통) & 3 & 6 & 9 & 12 & 15 \\
\textbf{가능성 2} (낮음) & 2 & 4 & 6 & 8 & 10 \\
\textbf{가능성 1} (희박) & 1 & 2 & 3 & 4 & 5 \\
\bottomrule
\end{tabularx}
\end{table}

\begin{table}[htbp]
\centering
\caption{리스크 등급 분류 기준}
\label{tab:risk_grade}
\begin{tabularx}{\textwidth}{p{2cm}p{2.5cm}X}
\toprule
\textbf{점수 범위} & \textbf{등급} & \textbf{대응 방침} \\
\midrule
20--25 & 극심(Critical) & 즉시 대응, 경영진 보고, 시스템 중단 검토 \\
12--19 & 높음(High) & 우선 대응, 30일 내 완화 계획 수립 \\
6--11 & 보통(Medium) & 계획적 대응, 분기 내 완화 조치 시행 \\
1--5 & 낮음(Low) & 수용 또는 모니터링, 연간 검토 \\
\bottomrule
\end{tabularx}
\end{table}

\subsection{리스크 평가 방법론: 4T 모델}
\label{sec:5.3.4}

식별되고 등급이 부여된 리스크에 대해 차등화된 대응 전략을 수립한다. 4T 모델은 네 가지 대응 유형을 제시한다.

\begin{enumerate}
 \item \textbf{처리(Treat)}: 리스크를 수용 가능한 수준으로 낮추기 위한 조치를 취한다. 기술적 완화(편향 제거, 암호화), 절차적 완화(검토 프로세스), 조직적 완화(교육, 역할 강화) 등을 포함한다.
 \item \textbf{전가(Transfer)}: 리스크를 제3자에게 전가한다. AI 보험 가입, 외부 서비스 계약에 책임 조항 포함, 공동 책임 체계 구축 등이 해당한다.
 \item \textbf{중단(Terminate)}: 리스크가 허용 수준을 초과하여 완화가 불가능한 경우 해당 AI 시스템 또는 활용 사례를 중단한다. 비용 대비 효과 분석을 통해 판단한다.
 \item \textbf{수용(Tolerate)}: 리스크가 허용 수준 이내이거나, 잔여 리스크를 경영진이 인지하고 수용하는 경우이다. 수용 결정은 반드시 문서화해야 한다.
\end{enumerate}

\subsection{리스크 평가 정량화 방법론}
\label{sec:5.3.5}

리스크 평가의 일관성과 객관성을 확보하기 위해 정량화된 점수 산정 기준을 수립한다.

\begin{table}[htbp]
\centering
\caption{리스크 발생 가능성 평가 기준}
\label{tab:likelihood_criteria}
\begin{tabularx}{\textwidth}{p{1cm}p{2cm}X}
\toprule
\textbf{점수} & \textbf{등급} & \textbf{판단 기준} \\
\midrule
5 & 확실 & 연 1회 이상 발생 예상, 유사 사례 다수 존재 \\
4 & 높음 & 2--3년 내 발생 예상, 업계 내 사례 존재 \\
3 & 보통 & 3--5년 내 발생 가능, 이론적 가능성 입증 \\
2 & 낮음 & 5--10년 내 발생 가능, 극히 제한적 사례 \\
1 & 희박 & 10년 이상 미발생 예상, 이론적 가능성만 존재 \\
\bottomrule
\end{tabularx}
\end{table}

\begin{table}[htbp]
\centering
\caption{리스크 영향도 평가 기준}
\label{tab:impact_criteria}
\begin{tabularx}{\textwidth}{p{1cm}p{2cm}X}
\toprule
\textbf{점수} & \textbf{등급} & \textbf{판단 기준} \\
\midrule
5 & 치명적 & 인명 피해, 대규모 법적 제재, 사업 존속 위협 \\
4 & 심각 & 중대 재무 손실, 규제 조치, 심각한 평판 손상 \\
3 & 보통 & 상당한 운영 차질, 고객 불만, 미디어 관심 \\
2 & 경미 & 일시적 서비스 저하, 제한적 영향 \\
1 & 미미 & 최소한의 영향, 일상적 운영 내 처리 가능 \\
\bottomrule
\end{tabularx}
\end{table}

\subsection{AI 리스크 자동 평가 엔진}
\label{sec:5.3.6}

NIST AI RMF의 4대 기능을 기반으로 한 자동 리스크 평가 엔진을 구현하면, 조직 내 다수의 AI 시스템에 대해 일관된 기준으로 리스크를 평가하고 추적할 수 있다.

\begin{lstlisting}[language=Python, caption={NIST AI RMF 기반 리스크 자동 평가 엔진}, label={lst:risk_engine}]
from dataclasses import dataclass, field
from typing import Dict, List, Tuple
from enum import Enum
import json


class RiskGrade(Enum):
    CRITICAL = "critical"
    HIGH = "high"
    MEDIUM = "medium"
    LOW = "low"


class ResponseStrategy(Enum):
    TREAT = "treat"
    TRANSFER = "transfer"
    TERMINATE = "terminate"
    TOLERATE = "tolerate"


@dataclass
class RiskItem:
    """개별 리스크 항목"""
    risk_id: str
    category: str       # technical, operational, compliance,
                        # strategic, ethical
    description: str
    likelihood: int      # 1-5
    impact: int          # 1-5
    owner: str
    strategy: str = ""
    mitigations: List[str] = field(default_factory=list)

    @property
    def score(self) -> int:
        return self.likelihood * self.impact

    @property
    def grade(self) -> RiskGrade:
        s = self.score
        if s >= 20:
            return RiskGrade.CRITICAL
        elif s >= 12:
            return RiskGrade.HIGH
        elif s >= 6:
            return RiskGrade.MEDIUM
        else:
            return RiskGrade.LOW


class NISTRiskAssessmentEngine:
    """NIST AI RMF 4대 기능 기반 리스크 평가 엔진"""

    def __init__(self, org_risk_appetite: int = 12):
        self.risk_appetite = org_risk_appetite
        self.risk_register: List[RiskItem] = []

    # -- Govern 기능 --
    def set_risk_appetite(self, level: int):
        """조직의 리스크 허용 수준 설정"""
        self.risk_appetite = level

    # -- Map 기능 --
    def identify_risks(
        self, system_id: str, system_info: Dict
    ) -> List[RiskItem]:
        """AI 시스템 맥락 기반 리스크 자동 식별"""
        risks = []

        # 기술 리스크 점검
        if system_info.get("model_type") == "deep_learning":
            risks.append(RiskItem(
                risk_id=f"{system_id}-T001",
                category="technical",
                description="딥러닝 모델 적대적 공격 취약성",
                likelihood=3, impact=4,
                owner=system_info.get("owner", ""),
            ))
        if system_info.get("input_type") == "user_text":
            risks.append(RiskItem(
                risk_id=f"{system_id}-T002",
                category="technical",
                description="프롬프트 인젝션 공격 위험",
                likelihood=4, impact=3,
                owner=system_info.get("owner", ""),
            ))

        # 윤리 리스크 점검
        if system_info.get("affects_individuals", False):
            risks.append(RiskItem(
                risk_id=f"{system_id}-E001",
                category="ethical",
                description="개인 대상 의사결정 편향 위험",
                likelihood=3, impact=4,
                owner=system_info.get("owner", ""),
            ))

        # 컴플라이언스 리스크 점검
        if system_info.get("processes_personal_data", False):
            risks.append(RiskItem(
                risk_id=f"{system_id}-C001",
                category="compliance",
                description="개인정보보호법 위반 위험",
                likelihood=2, impact=5,
                owner=system_info.get("owner", ""),
            ))

        # 운영 리스크 점검
        if system_info.get("is_production", False):
            risks.append(RiskItem(
                risk_id=f"{system_id}-O001",
                category="operational",
                description="데이터 드리프트로 인한 성능 저하",
                likelihood=4, impact=3,
                owner=system_info.get("owner", ""),
            ))

        return risks

    # -- Measure 기능 --
    def evaluate_risks(
        self, risks: List[RiskItem]
    ) -> Dict[str, object]:
        """리스크 정량 평가 및 등급 산정"""
        results = {
            "total_risks": len(risks),
            "by_grade": {g.value: 0 for g in RiskGrade},
            "by_category": {},
            "above_appetite": [],
            "risk_details": [],
        }

        for risk in risks:
            grade = risk.grade.value
            results["by_grade"][grade] += 1

            cat = risk.category
            if cat not in results["by_category"]:
                results["by_category"][cat] = []
            results["by_category"][cat].append(risk.risk_id)

            if risk.score > self.risk_appetite:
                results["above_appetite"].append({
                    "id": risk.risk_id,
                    "score": risk.score,
                    "grade": grade,
                })

            results["risk_details"].append({
                "id": risk.risk_id,
                "score": risk.score,
                "grade": grade,
                "category": cat,
            })

        return results

    # -- Manage 기능 --
    def recommend_strategy(
        self, risk: RiskItem
    ) -> ResponseStrategy:
        """4T 모델 기반 대응 전략 추천"""
        if risk.score >= 20:
            return ResponseStrategy.TERMINATE
        elif risk.score >= 12:
            return ResponseStrategy.TREAT
        elif risk.score >= 6:
            if risk.category == "compliance":
                return ResponseStrategy.TREAT
            return ResponseStrategy.TRANSFER
        else:
            return ResponseStrategy.TOLERATE

    def generate_report(
        self, system_id: str, risks: List[RiskItem]
    ) -> str:
        """리스크 평가 보고서 생성"""
        evaluation = self.evaluate_risks(risks)
        lines = [
            f"=== 리스크 평가 보고서: {system_id} ===",
            f"총 리스크 수: {evaluation['total_risks']}",
            f"등급별 분포: {evaluation['by_grade']}",
            f"허용 수준 초과: "
            f"{len(evaluation['above_appetite'])}건",
            "",
        ]
        for risk in risks:
            strategy = self.recommend_strategy(risk)
            lines.append(
                f"[{risk.risk_id}] 점수={risk.score} "
                f"등급={risk.grade.value} "
                f"전략={strategy.value}"
            )
        return "\n".join(lines)


# 사용 예시
if __name__ == "__main__":
    engine = NISTRiskAssessmentEngine(org_risk_appetite=12)

    system_info = {
        "model_type": "deep_learning",
        "input_type": "user_text",
        "affects_individuals": True,
        "processes_personal_data": True,
        "is_production": True,
        "owner": "AI개발팀",
    }

    risks = engine.identify_risks("SYS-001", system_info)
    report = engine.generate_report("SYS-001", risks)
    print(report)
\end{lstlisting}

\begin{practicaltool}{AI 리스크 레지스터 운영 가이드}
AI 리스크 레지스터는 조직 내 모든 AI 리스크를 중앙에서 관리하는 핵심 도구이다. 효과적인 운영을 위해 다음 원칙을 준수한다.

\textbf{표: AI 리스크 레지스터 필수 항목 (확장)}
\label{tab:risk_register_extended}

\centering
\begin{tabularx}{\textwidth}{p{2.8cm}X}
\toprule
\textbf{항목} & \textbf{설명} \\
\midrule
리스크 ID & 고유 식별자 (예: R-AI-2026-001) \\
등록일 & 리스크 최초 식별 일자 \\
카테고리 & 기술, 운영, 윤리, 컴플라이언스, 전략 중 분류 \\
설명 & 리스크에 대한 구체적 기술 \\
관련 시스템 & 영향받는 AI 시스템 목록 \\
발생 가능성 & 1--5점 척도 평가 \\
영향도 & 1--5점 척도 평가 \\
리스크 점수 & 가능성 $\times$ 영향도 \\
리스크 등급 & 극심, 높음, 보통, 낮음 \\
대응 전략 & 4T 중 선택 (처리, 전가, 중단, 수용) \\
완화 조치 & 구체적 완화 활동 기술 \\
책임자 & 리스크 소유자(Risk Owner) \\
완화 기한 & 완화 조치 완료 목표일 \\
상태 & 신규, 진행 중, 완화 완료, 수용, 종결 \\
잔여 리스크 & 완화 후 잔여 리스크 수준 \\
최종 검토일 & 마지막 검토 수행일 \\
\bottomrule
\end{tabularx}

\textbf{운영 원칙:}
\begin{enumerate}
 \item 모든 AI 시스템은 도입 시 리스크 레지스터에 등록한다.
 \item 리스크 점수 12점 이상은 매월 검토하고, 6점 이상은 분기별 검토한다.
 \item 리스크 등급 변동 시 즉시 이해관계자에게 통보한다.
 \item 완화 조치 완료 시 잔여 리스크를 재평가한다.
 \item 리스크 레지스터는 분기별 경영진 보고에 포함한다.
\end{enumerate}
\end{practicaltool}

%----------------------------------------------------------
\section{정책 준수 모니터링 체계}
\label{sec:5.4}

정책이 수립되었다 하더라도 이를 지속적으로 모니터링하지 않으면 형식적 문서에 그칠 위험이 크다. 정책 준수 모니터링 체계는 조직 내 AI 활동이 수립된 정책을 실제로 따르고 있는지를 체계적으로 추적하고 검증하는 메커니즘이다.

\subsection{KPI 기반 정책 준수 추적}
\label{sec:5.4.1}

정책 준수 상태를 객관적으로 측정하기 위해 핵심 성과 지표(KPI, Key Performance Indicator)를 정의하고 추적한다.

\begin{table}[htbp]
\centering
\caption{AI 거버넌스 정책 준수 KPI}
\label{tab:compliance_kpi}
\begin{tabularx}{\textwidth}{p{3.5cm}p{2cm}p{2cm}X}
\toprule
\textbf{KPI} & \textbf{목표} & \textbf{측정 주기} & \textbf{산출 방법} \\
\midrule
영향평가 완료율 & 100\% & 월간 & 고위험 시스템 중 영향평가 완료 비율 \\
모델 카드 작성률 & $\geq$95\% & 월간 & 배포 모델 중 모델 카드 보유 비율 \\
공정성 기준 충족률 & $\geq$90\% & 월간 & 공정성 메트릭 임계값 충족 비율 \\
인시던트 대응 시간 & $\leq$4시간 & 건별 & 인시던트 보고~초기 대응 소요 시간 \\
정책 교육 이수율 & $\geq$95\% & 분기 & 대상 인원 중 교육 이수 비율 \\
리스크 레지스터 갱신율 & 100\% & 월간 & 기한 내 검토 완료된 리스크 비율 \\
배포 게이트 통과율 & 추적 & 월간 & 배포 요청 대비 승인 비율 \\
데이터 품질 점수 & $\geq$0.85 & 월간 & 데이터 완전성, 정확성, 적시성 평균 \\
\bottomrule
\end{tabularx}
\end{table}

\subsection{정책 위반 탐지 및 에스컬레이션}
\label{sec:5.4.2}

정책 위반이 발생했을 때 신속하게 탐지하고 적절한 수준으로 보고하는 에스컬레이션 프로세스가 필요하다.

\begin{table}[htbp]
\centering
\caption{정책 위반 심각도 분류 및 에스컬레이션 경로}
\label{tab:violation_escalation}
\begin{tabularx}{\textwidth}{p{1.5cm}p{3cm}p{2.5cm}X}
\toprule
\textbf{등급} & \textbf{위반 유형} & \textbf{보고 대상} & \textbf{대응 시한} \\
\midrule
심각 & 법규 위반, 개인정보 유출 & CEO, 이사회, 규제 기관 & 즉시 (1시간 이내) \\
높음 & 고위험 정책 미준수 & CDO/CTO, 거버넌스 위원회 & 4시간 이내 \\
보통 & 표준/지침 미준수 & 부서장, 거버넌스팀 & 24시간 이내 \\
낮음 & 절차 미준수 & 팀장, 담당자 & 1주일 이내 \\
\bottomrule
\end{tabularx}
\end{table}

에스컬레이션 프로세스는 다음 원칙을 따른다:
\begin{enumerate}
 \item \textbf{자동 탐지 우선}: 가능한 한 자동화된 모니터링 시스템에서 위반을 탐지한다.
 \item \textbf{무처벌 보고 문화}: 위반 보고자를 보호하여 자발적 보고를 장려한다.
 \item \textbf{근본 원인 분석}: 위반 처리 후 반드시 근본 원인 분석(Root Cause Analysis)을 수행한다.
 \item \textbf{재발 방지}: 동일 유형의 위반이 반복되지 않도록 시스템적 개선 조치를 수립한다.
\end{enumerate}

\subsection{정기 감사와 임시 감사}
\label{sec:5.4.3}

정책 준수를 체계적으로 검증하기 위해 정기 감사와 임시 감사를 병행한다. 감사 체계의 상세한 내용은 제11장(\ref{ch:audit_certification}장)에서 다루며, 여기서는 정책 관점에서의 핵심 사항을 정리한다.

\begin{itemize}
 \item \textbf{정기 감사}: 연 1회 이상 전체 AI 시스템에 대한 정책 준수 감사를 수행한다. 감사 범위는 정책 문서 현행화, 리스크 레지스터 관리, 영향평가 수행 현황, KPI 달성도를 포함한다.
 \item \textbf{임시 감사}: 중대 인시던트 발생, 규제 환경 변화, 신규 고위험 시스템 도입 시 임시 감사를 시행한다.
 \item \textbf{자가 점검}: 각 AI 시스템 담당팀이 분기별로 자가 점검을 수행하고 결과를 거버넌스팀에 보고한다.
\end{itemize}

\subsection{정책 준수 대시보드 설계}
\label{sec:5.4.4}

정책 준수 현황을 경영진과 실무자 모두가 한눈에 파악할 수 있도록 대시보드를 설계한다.

\begin{table}[htbp]
\centering
\caption{정책 준수 대시보드 구성 요소}
\label{tab:dashboard_components}
\begin{tabularx}{\textwidth}{p{3cm}p{2.5cm}X}
\toprule
\textbf{대시보드 영역} & \textbf{시각화 유형} & \textbf{표시 내용} \\
\midrule
전사 준수 현황 & 게이지 차트 & 전체 정책 준수율, 목표 대비 현황 \\
리스크 히트맵 & 히트맵 & 시스템별 리스크 점수 분포 \\
KPI 추이 & 라인 차트 & 주요 KPI의 월별 추이 \\
위반 현황 & 막대 차트 & 등급별, 카테고리별 위반 건수 \\
시스템별 상태 & 테이블 & 개별 AI 시스템의 준수 상태 요약 \\
기한 관리 & 캘린더 & 영향평가, 감사, 검토 일정 \\
\bottomrule
\end{tabularx}
\end{table}

%----------------------------------------------------------
\section{Policy-as-Code 자동화 체계}
\label{sec:5.5}

AI 기술의 빠른 변화 속도에 맞춰 정책을 시스템에 내재화하는 Policy-as-Code(PaC) 개념이 강조된다. PaC는 인간이 읽을 수 있는 정책 문서를 기계가 실행할 수 있는 코드로 변환하여, 정책 준수를 자동으로 검증하고 강제하는 패러다임이다.

\subsection{OPA 기반 거버넌스 자동화}
\label{sec:5.5.1}

OPA(Open Policy Agent)는 정책을 코드로 표현하고 의사결정을 자동화하는 오픈소스 엔진이다. AI 거버넌스에서 OPA를 활용하면 모델 배포, 데이터 접근, 서비스 구성 등에 대한 정책을 통일된 방식으로 관리할 수 있다.

OPA의 핵심 구성 요소는 다음과 같다:
\begin{itemize}
 \item \textbf{정책(Policy)}: Rego 언어로 작성된 규칙 집합
 \item \textbf{데이터(Data)}: 정책 결정에 필요한 맥락 정보(JSON)
 \item \textbf{질의(Query)}: 정책 결정을 요청하는 입력
 \item \textbf{결정(Decision)}: 정책 평가 결과(허용/거부/조건부 허용)
\end{itemize}

\subsection{CI/CD 파이프라인 통합}
\label{sec:5.5.2}

Policy-as-Code를 CI/CD(지속적 통합/지속적 배포) 파이프라인에 통합하면, 모델이 배포되기 전에 자동으로 정책 준수를 검증할 수 있다. 이를 ``배포 게이트(Deployment Gate)''라 한다.

배포 게이트의 검증 항목은 다음과 같다:
\begin{enumerate}
 \item \textbf{성능 기준}: 정확도, F1 점수, AUC 등이 임계값 이상인지 확인
 \item \textbf{공정성 기준}: 인구통계학적 패리티 비율(Demographic Parity Ratio)이 기준 이상인지 확인
 \item \textbf{프라이버시 기준}: 차분 프라이버시 엡실론 값, 비식별화 수준 확인
 \item \textbf{문서 기준}: 모델 카드, 데이터시트, 영향평가 보고서 존재 여부 확인
 \item \textbf{보안 기준}: 취약점 스캔 결과, 접근 통제 설정 확인
 \item \textbf{승인 기준}: 필요한 승인이 모두 완료되었는지 확인
\end{enumerate}

\subsection{Policy-as-Code 종합 구현}
\label{sec:5.5.3}

다음은 배포 게이트, 공정성 검증, 프라이버시 검증을 포함한 종합적인 Policy-as-Code 구현이다.

\begin{lstlisting}[language=Python, caption={Policy-as-Code 종합 구현: 배포 게이트}, label={lst:pac_full}]
from dataclasses import dataclass, field
from typing import Dict, List, Optional
from enum import Enum
from datetime import datetime
import json


class GateDecision(Enum):
    APPROVED = "approved"
    REJECTED = "rejected"
    CONDITIONAL = "conditional"


@dataclass
class PolicyRule:
    """개별 정책 규칙"""
    rule_id: str
    name: str
    category: str
    threshold: float
    operator: str  # gte, lte, eq, exists
    description: str


@dataclass
class CheckResult:
    """개별 검증 결과"""
    rule_id: str
    passed: bool
    actual_value: object
    expected: str
    message: str


@dataclass
class GateResult:
    """배포 게이트 종합 결과"""
    model_id: str
    decision: GateDecision
    checks: List[CheckResult]
    timestamp: str = field(
        default_factory=lambda: datetime.now().isoformat()
    )

    @property
    def pass_rate(self) -> float:
        if not self.checks:
            return 0.0
        passed = sum(1 for c in self.checks if c.passed)
        return round(passed / len(self.checks), 3)


class AIDeploymentGate:
    """AI 모델 배포 게이트 - Policy-as-Code 엔진"""

    def __init__(self):
        self.rules: List[PolicyRule] = []
        self._load_default_rules()

    def _load_default_rules(self):
        """기본 정책 규칙 로드"""
        self.rules = [
            # 성능 기준
            PolicyRule(
                "PERF-001", "최소 정확도",
                "performance", 0.85, "gte",
                "모델 정확도 85% 이상"
            ),
            PolicyRule(
                "PERF-002", "최소 F1 점수",
                "performance", 0.80, "gte",
                "F1 점수 80% 이상"
            ),
            # 공정성 기준
            PolicyRule(
                "FAIR-001", "인구통계 패리티 비율",
                "fairness", 0.80, "gte",
                "DP 비율 0.80 이상"
            ),
            PolicyRule(
                "FAIR-002", "균등 기회 차이",
                "fairness", 0.10, "lte",
                "균등 기회 차이 0.10 이하"
            ),
            # 프라이버시 기준
            PolicyRule(
                "PRIV-001", "차분 프라이버시 엡실론",
                "privacy", 1.0, "lte",
                "DP 엡실론 1.0 이하"
            ),
            PolicyRule(
                "PRIV-002", "비식별화 적용",
                "privacy", 1.0, "eq",
                "비식별화 기법 적용 여부"
            ),
            # 문서 기준
            PolicyRule(
                "DOC-001", "모델 카드 존재",
                "documentation", 1.0, "exists",
                "모델 카드 작성 완료"
            ),
            PolicyRule(
                "DOC-002", "영향평가 완료",
                "documentation", 1.0, "exists",
                "영향평가 보고서 존재"
            ),
        ]

    def _evaluate_rule(
        self, rule: PolicyRule, value: object
    ) -> CheckResult:
        """개별 규칙 평가"""
        passed = False
        if rule.operator == "gte":
            passed = float(value) >= rule.threshold
        elif rule.operator == "lte":
            passed = float(value) <= rule.threshold
        elif rule.operator == "eq":
            passed = float(value) == rule.threshold
        elif rule.operator == "exists":
            passed = value is not None and value != ""

        expected = (
            f"{rule.operator} {rule.threshold}"
        )
        return CheckResult(
            rule_id=rule.rule_id,
            passed=passed,
            actual_value=value,
            expected=expected,
            message=rule.description,
        )

    def evaluate(
        self, model_id: str, metrics: Dict
    ) -> GateResult:
        """배포 게이트 종합 평가"""
        checks = []
        metric_map = {
            "PERF-001": metrics.get("accuracy"),
            "PERF-002": metrics.get("f1_score"),
            "FAIR-001": metrics.get("dp_ratio"),
            "FAIR-002": metrics.get("eq_opportunity_diff"),
            "PRIV-001": metrics.get("dp_epsilon"),
            "PRIV-002": metrics.get("anonymized"),
            "DOC-001": metrics.get("model_card"),
            "DOC-002": metrics.get("impact_assessment"),
        }

        for rule in self.rules:
            value = metric_map.get(rule.rule_id)
            result = self._evaluate_rule(rule, value)
            checks.append(result)

        # 의사결정 로직
        failed = [c for c in checks if not c.passed]
        critical_fail = any(
            c.rule_id.startswith(("FAIR", "PRIV"))
            for c in failed
        )

        if not failed:
            decision = GateDecision.APPROVED
        elif critical_fail:
            decision = GateDecision.REJECTED
        else:
            decision = GateDecision.CONDITIONAL

        return GateResult(
            model_id=model_id,
            decision=decision,
            checks=checks,
        )

    def generate_report(self, result: GateResult) -> str:
        """배포 게이트 결과 보고서 생성"""
        lines = [
            "=" * 50,
            f"배포 게이트 결과: {result.model_id}",
            f"판정: {result.decision.value}",
            f"통과율: {result.pass_rate * 100:.1f}%",
            f"평가 시점: {result.timestamp}",
            "-" * 50,
        ]
        for c in result.checks:
            status = "PASS" if c.passed else "FAIL"
            lines.append(
                f"[{status}] {c.rule_id}: "
                f"{c.message} "
                f"(실제: {c.actual_value}, "
                f"기준: {c.expected})"
            )
        lines.append("=" * 50)
        return "\n".join(lines)


# 사용 예시
if __name__ == "__main__":
    gate = AIDeploymentGate()

    metrics = {
        "accuracy": 0.92,
        "f1_score": 0.88,
        "dp_ratio": 0.78,
        "eq_opportunity_diff": 0.05,
        "dp_epsilon": 0.8,
        "anonymized": 1.0,
        "model_card": "model_card_v2.md",
        "impact_assessment": "ia_report_2026.pdf",
    }

    result = gate.evaluate("MODEL-2026-042", metrics)
    print(gate.generate_report(result))
\end{lstlisting}

%----------------------------------------------------------
\section{인시던트 관리 및 위기 대응}
\label{sec:5.6}

AI 시스템 운영 중 발생하는 인시던트(Incident)에 대한 체계적인 관리와 신속한 대응은 AI 거버넌스의 핵심 요소이다. AI 인시던트는 전통적 IT 인시던트와 달리 편향에 의한 차별, 환각(Hallucination)에 의한 오정보 생성, 데이터 드리프트에 의한 점진적 성능 저하 등 AI 고유의 특성을 반영해야 한다.

\subsection{AI 인시던트 분류 체계}
\label{sec:5.6.1}

AI 인시던트의 심각도를 체계적으로 분류하여 대응 수준을 결정한다.

\begin{table}[htbp]
\centering
\caption{AI 인시던트 심각도 분류}
\label{tab:incident_severity}
\begin{tabularx}{\textwidth}{p{1.5cm}p{2cm}X p{2.5cm}}
\toprule
\textbf{등급} & \textbf{심각도} & \textbf{기준} & \textbf{대응 시한} \\
\midrule
1등급 & 치명적 & 인명 피해 발생 또는 위험, 대규모 개인정보 유출, 법적 제재 발생 & 즉시 (15분 이내) \\
\midrule
2등급 & 심각 & 서비스 전면 중단, 차별적 결과 대량 발생, 규제 기관 통보 필요 & 1시간 이내 \\
\midrule
3등급 & 보통 & 부분적 서비스 장애, 성능 저하, 제한적 범위의 편향 발견 & 4시간 이내 \\
\midrule
4등급 & 경미 & 일시적 오류, 단일 건 민원, 경미한 성능 변동 & 24시간 이내 \\
\bottomrule
\end{tabularx}
\end{table}

\begin{table}[htbp]
\centering
\caption{AI 고유 인시던트 유형 예시}
\label{tab:ai_incident_types}
\begin{tabularx}{\textwidth}{p{3cm}p{2cm}X}
\toprule
\textbf{인시던트 유형} & \textbf{관련 리스크} & \textbf{사례} \\
\midrule
모델 성능 급락 & 기술 & 데이터 드리프트로 예측 정확도 30\% 이상 하락 \\
편향 탐지 & 윤리 & 특정 인구 집단에 대한 거부율이 타 집단 대비 2배 이상 \\
프롬프트 인젝션 & 보안 & 악의적 입력으로 시스템 보호 장치 우회 \\
환각 발생 & 기술 & 생성형 AI가 허위 정보를 사실처럼 출력 \\
개인정보 유출 & 컴플라이언스 & 모델이 학습 데이터의 개인정보를 출력 \\
적대적 공격 & 보안 & 조작된 입력으로 모델 오작동 유발 \\
\bottomrule
\end{tabularx}
\end{table}

\subsection{인시던트 대응 프로세스}
\label{sec:5.6.2}

NIST 사이버보안 프레임워크의 인시던트 대응 체계를 AI 맥락에 맞게 조정한 5단계 프로세스를 적용한다.

\paragraph{1단계: 탐지 및 보고(Detection \& Reporting)}
AI 모니터링 시스템의 자동 경보, 사용자 민원, 내부 보고를 통해 인시던트를 탐지한다. 모든 인시던트는 발생 즉시 인시던트 관리 시스템에 등록한다.

\paragraph{2단계: 분류 및 에스컬레이션(Triage \& Escalation)}
인시던트의 심각도를 평가하고 등급을 부여한다. 등급에 따라 적절한 대응팀을 소집하고 이해관계자에게 통보한다.

\paragraph{3단계: 봉쇄 및 완화(Containment \& Mitigation)}
인시던트의 확산을 방지하기 위한 즉각적 조치를 취한다. AI 시스템 일시 중단, 이전 버전 롤백, 인적 감독 강화, 영향받는 사용자 통보 등을 포함한다.

\paragraph{4단계: 근본 원인 분석(Root Cause Analysis)}
인시던트의 근본 원인을 파악한다. 데이터 품질 문제, 모델 설계 결함, 운영 프로세스 미비, 외부 공격 등 원인을 체계적으로 분석한다.

\paragraph{5단계: 복구 및 학습(Recovery \& Learning)}
시스템을 정상 상태로 복구하고, 인시던트로부터 학습한 교훈을 정책과 프로세스에 반영한다. 사후 검토(Post-mortem) 보고서를 작성하고 조직 전체에 공유한다.

\subsection{인시던트 보고 및 학습 체계}
\label{sec:5.6.3}

인시던트로부터의 조직 학습은 동일한 사고의 재발을 방지하는 핵심 메커니즘이다.

\begin{itemize}
 \item \textbf{사후 검토(Post-mortem)}: 모든 2등급 이상 인시던트에 대해 사후 검토를 수행한다. 비난하지 않는 문화(Blameless Culture)를 기반으로 한다.
 \item \textbf{인시던트 데이터베이스}: 모든 인시던트를 데이터베이스에 기록하여 패턴을 분석하고, 예방적 조치를 도출한다.
 \item \textbf{교훈 공유}: 인시던트 교훈을 조직 전체에 공유하여, 유사 시스템에 대한 사전 점검을 촉진한다.
 \item \textbf{정책 피드백}: 인시던트 교훈을 정책 및 표준 개정에 반영하여 거버넌스 체계를 지속적으로 강화한다.
\end{itemize}

\subsection{위기 커뮤니케이션 계획}
\label{sec:5.6.4}

AI 인시던트가 대외적으로 확산될 경우를 대비한 위기 커뮤니케이션 계획을 사전에 수립해야 한다.

\begin{table}[htbp]
\centering
\caption{위기 커뮤니케이션 대상 및 내용}
\label{tab:crisis_communication}
\begin{tabularx}{\textwidth}{p{2.5cm}p{2cm}X}
\toprule
\textbf{커뮤니케이션 대상} & \textbf{시한} & \textbf{전달 내용} \\
\midrule
경영진/이사회 & 1시간 이내 & 인시던트 개요, 영향 범위, 초기 대응 현황 \\
규제 기관 & 법정 기한 내 & 법적 보고 양식에 따른 공식 보고 \\
영향받는 이용자 & 24시간 이내 & 사고 사실, 영향 범위, 보호 조치, 구제 방법 \\
미디어 & 48시간 이내 & 공식 성명, 대응 계획, 재발 방지 약속 \\
내부 직원 & 4시간 이내 & 사고 개요, 행동 지침, 외부 문의 대응 방법 \\
\bottomrule
\end{tabularx}
\end{table}

\subsection{인시던트 관리 시스템 구현}
\label{sec:5.6.5}

인시던트 관리의 효율성을 높이기 위해 자동화된 시스템을 구축한다. 다음은 인시던트의 등록, 분류, 추적, 보고를 통합 관리하는 파이썬 구현이다.

\begin{lstlisting}[language=Python, caption={AI 인시던트 관리 시스템}, label={lst:incident_mgmt}]
from dataclasses import dataclass, field
from typing import Dict, List, Optional
from enum import Enum
from datetime import datetime


class Severity(Enum):
    CRITICAL = 1    # 치명적
    MAJOR = 2       # 심각
    MODERATE = 3    # 보통
    MINOR = 4       # 경미


class IncidentStatus(Enum):
    REPORTED = "reported"
    TRIAGED = "triaged"
    CONTAINED = "contained"
    INVESTIGATING = "investigating"
    RESOLVED = "resolved"
    CLOSED = "closed"


class IncidentType(Enum):
    PERFORMANCE_DEGRADATION = "performance_degradation"
    BIAS_DETECTED = "bias_detected"
    PROMPT_INJECTION = "prompt_injection"
    HALLUCINATION = "hallucination"
    DATA_BREACH = "data_breach"
    ADVERSARIAL_ATTACK = "adversarial_attack"
    DATA_DRIFT = "data_drift"
    SERVICE_OUTAGE = "service_outage"


@dataclass
class Incident:
    """AI 인시던트"""
    incident_id: str
    title: str
    incident_type: IncidentType
    severity: Severity
    system_id: str
    description: str
    reported_by: str
    reported_at: str = field(
        default_factory=lambda: datetime.now().isoformat()
    )
    status: IncidentStatus = IncidentStatus.REPORTED
    assigned_to: str = ""
    timeline: List[Dict] = field(default_factory=list)
    root_cause: str = ""
    resolution: str = ""
    lessons_learned: List[str] = field(
        default_factory=list
    )

    def add_timeline_entry(
        self, action: str, actor: str
    ):
        self.timeline.append({
            "timestamp": datetime.now().isoformat(),
            "action": action,
            "actor": actor,
        })


# 에스컬레이션 규칙
ESCALATION_RULES = {
    Severity.CRITICAL: {
        "response_time_minutes": 15,
        "notify": [
            "ceo", "cto", "legal",
            "governance_committee"
        ],
        "actions": [
            "시스템 즉시 중단",
            "인시던트 대응팀 긴급 소집",
            "규제 기관 통보 준비",
        ],
    },
    Severity.MAJOR: {
        "response_time_minutes": 60,
        "notify": ["cto", "governance_team"],
        "actions": [
            "인적 감독 전환",
            "영향 범위 긴급 파악",
            "사용자 고지 준비",
        ],
    },
    Severity.MODERATE: {
        "response_time_minutes": 240,
        "notify": ["team_lead", "governance_team"],
        "actions": [
            "모니터링 강화",
            "원인 분석 착수",
            "임시 완화 조치",
        ],
    },
    Severity.MINOR: {
        "response_time_minutes": 1440,
        "notify": ["team_lead"],
        "actions": [
            "일반 대응 절차",
            "다음 정기 검토에 포함",
        ],
    },
}


class AIIncidentManager:
    """AI 인시던트 관리 시스템"""

    def __init__(self):
        self.incidents: Dict[str, Incident] = {}
        self._counter = 0

    def report_incident(
        self,
        title: str,
        incident_type: IncidentType,
        severity: Severity,
        system_id: str,
        description: str,
        reported_by: str,
    ) -> Incident:
        """인시던트 등록"""
        self._counter += 1
        incident_id = (
            f"INC-{datetime.now().year}-"
            f"{self._counter:04d}"
        )

        incident = Incident(
            incident_id=incident_id,
            title=title,
            incident_type=incident_type,
            severity=severity,
            system_id=system_id,
            description=description,
            reported_by=reported_by,
        )
        incident.add_timeline_entry(
            "인시던트 등록", reported_by
        )

        self.incidents[incident_id] = incident

        # 자동 에스컬레이션
        self._escalate(incident)

        return incident

    def _escalate(self, incident: Incident):
        """에스컬레이션 규칙 적용"""
        rules = ESCALATION_RULES[incident.severity]
        incident.add_timeline_entry(
            f"에스컬레이션: "
            f"{rules['notify']}에 통보, "
            f"대응 시한 "
            f"{rules['response_time_minutes']}분",
            "system",
        )
        for action in rules["actions"]:
            incident.add_timeline_entry(
                f"권고 조치: {action}", "system"
            )

    def triage(
        self,
        incident_id: str,
        assigned_to: str,
        updated_severity: Optional[Severity] = None,
    ):
        """인시던트 분류"""
        inc = self.incidents[incident_id]
        inc.status = IncidentStatus.TRIAGED
        inc.assigned_to = assigned_to
        if updated_severity:
            inc.severity = updated_severity
            self._escalate(inc)
        inc.add_timeline_entry(
            f"{assigned_to}에게 할당", "triage_team"
        )

    def contain(
        self, incident_id: str, actions: List[str]
    ):
        """봉쇄 조치"""
        inc = self.incidents[incident_id]
        inc.status = IncidentStatus.CONTAINED
        for action in actions:
            inc.add_timeline_entry(
                f"봉쇄 조치: {action}", inc.assigned_to
            )

    def resolve(
        self,
        incident_id: str,
        root_cause: str,
        resolution: str,
        lessons: List[str],
    ):
        """인시던트 해결"""
        inc = self.incidents[incident_id]
        inc.status = IncidentStatus.RESOLVED
        inc.root_cause = root_cause
        inc.resolution = resolution
        inc.lessons_learned = lessons
        inc.add_timeline_entry(
            "인시던트 해결 완료", inc.assigned_to
        )

    def generate_postmortem(
        self, incident_id: str
    ) -> str:
        """사후 검토 보고서 생성"""
        inc = self.incidents[incident_id]
        lines = [
            "=" * 50,
            f"사후 검토 보고서: {inc.incident_id}",
            "=" * 50,
            f"제목: {inc.title}",
            f"유형: {inc.incident_type.value}",
            f"심각도: {inc.severity.name}",
            f"시스템: {inc.system_id}",
            f"보고자: {inc.reported_by}",
            f"보고 시각: {inc.reported_at}",
            f"상태: {inc.status.value}",
            "",
            "--- 근본 원인 ---",
            inc.root_cause or "분석 중",
            "",
            "--- 해결 조치 ---",
            inc.resolution or "진행 중",
            "",
            "--- 교훈 ---",
        ]
        for lesson in inc.lessons_learned:
            lines.append(f"  - {lesson}")
        lines.append("")
        lines.append("--- 타임라인 ---")
        for entry in inc.timeline:
            lines.append(
                f"  [{entry['timestamp']}] "
                f"{entry['action']} "
                f"(by {entry['actor']})"
            )
        return "\n".join(lines)

    def get_dashboard_data(self) -> Dict:
        """대시보드용 데이터 산출"""
        data = {
            "total": len(self.incidents),
            "by_severity": {},
            "by_status": {},
            "by_type": {},
            "open_incidents": [],
        }
        for inc in self.incidents.values():
            sev = inc.severity.name
            data["by_severity"][sev] = (
                data["by_severity"].get(sev, 0) + 1
            )
            st = inc.status.value
            data["by_status"][st] = (
                data["by_status"].get(st, 0) + 1
            )
            tp = inc.incident_type.value
            data["by_type"][tp] = (
                data["by_type"].get(tp, 0) + 1
            )
            if inc.status not in (
                IncidentStatus.RESOLVED,
                IncidentStatus.CLOSED
            ):
                data["open_incidents"].append(
                    inc.incident_id
                )
        return data


# 사용 예시
if __name__ == "__main__":
    manager = AIIncidentManager()

    # 인시던트 등록
    inc = manager.report_incident(
        title="채용 AI 성별 편향 탐지",
        incident_type=IncidentType.BIAS_DETECTED,
        severity=Severity.MAJOR,
        system_id="SYS-RECRUIT-001",
        description=(
            "채용 심사 AI에서 여성 지원자 합격률이 "
            "남성 대비 65% 수준으로 하락"
        ),
        reported_by="monitoring_system",
    )

    # 분류
    manager.triage(inc.incident_id, "AI거버넌스팀")

    # 봉쇄
    manager.contain(inc.incident_id, [
        "해당 AI 시스템 일시 중단",
        "수동 심사로 전환",
        "영향 받은 지원자 목록 확보",
    ])

    # 해결
    manager.resolve(
        inc.incident_id,
        root_cause=(
            "학습 데이터 내 역사적 채용 편향 반영"
        ),
        resolution=(
            "편향 완화 기법 적용 후 재학습, "
            "공정성 메트릭 임계값 강화"
        ),
        lessons=[
            "학습 데이터 편향 사전 점검 절차 추가",
            "공정성 모니터링 임계값 조정",
            "분기별 편향 감사 의무화",
        ],
    )

    print(manager.generate_postmortem(inc.incident_id))
    print("\n대시보드:", manager.get_dashboard_data())
\end{lstlisting}

\begin{practicaltool}{AI 인시던트 대응 플레이북}
인시던트 유형별로 사전 정의된 대응 절차를 정리한 플레이북이다. 실제 인시던트 발생 시 신속한 대응을 위해 활용한다.

\textbf{유형 1: 편향/차별 탐지}
\begin{enumerate}
 \item 해당 AI 시스템의 자동 의사결정 기능 즉시 중단
 \item 인적 감독 모드로 전환
 \item 영향받은 이용자 범위 파악
 \item 공정성 메트릭 상세 분석 수행
 \item 편향 원인 분석(데이터, 모델, 후처리 단계별)
 \item 완화 조치 적용 후 검증
 \item 영향받은 이용자에게 고지 및 구제 조치
\end{enumerate}

\textbf{유형 2: 개인정보 유출}
\begin{enumerate}
 \item 해당 시스템 접근 즉시 차단
 \item 개인정보보호 책임자(DPO) 및 법무팀 즉시 통보
 \item 유출 범위 및 영향 정보주체 파악
 \item 72시간 이내 개인정보보호위원회 신고(GDPR/한국법)
 \item 정보주체에게 유출 사실 통지
 \item 기술적 원인 분석 및 보안 조치 강화
 \item 사후 검토 및 재발 방지 대책 수립
\end{enumerate}

\textbf{유형 3: 모델 성능 급락}
\begin{enumerate}
 \item 성능 저하 범위 및 심각도 파악
 \item 데이터 드리프트 여부 점검
 \item 이전 안정 버전으로 롤백 결정
 \item 원인 분석(데이터 변화, 인프라 이상, 외부 공격)
 \item 모델 재학습 또는 파이프라인 수정
 \item 강화된 모니터링 하에 재배포
\end{enumerate}

\textbf{유형 4: 적대적 공격/프롬프트 인젝션}
\begin{enumerate}
 \item 공격 패턴 즉시 식별 및 차단
 \item 보안팀 통보 및 합동 대응
 \item 입력 필터링 규칙 긴급 업데이트
 \item 피해 범위 파악 및 오염된 출력 회수
 \item 공격 벡터 분석 및 방어 메커니즘 강화
 \item 유사 시스템 대상 사전 점검
\end{enumerate}
\end{practicaltool}

%----------------------------------------------------------
\section*{제5장 요약}

본 장에서는 AI 정책 및 리스크 관리를 위한 포괄적 체계를 다루었다. 주요 내용을 정리하면 다음과 같다.

\begin{enumerate}
 \item \textbf{정책 아키텍처}: 3계층(정책--표준/지침--절차/매뉴얼) 구조를 통해 전략적 방향부터 실무 지침까지 체계적으로 관리한다. AI 윤리 정책, AI 활용 정책, AI 보안 정책이 핵심 정책 문서를 구성한다.

 \item \textbf{AI 영향평가}: EU AI Act와 한국 AI 기본법의 요구사항을 반영한 4단계 영향평가 절차(사전 스크리닝, 상세 평가, 완화 조치, 모니터링)를 수립하고, 자동화 도구를 통해 일관된 평가를 수행한다.

 \item \textbf{리스크 분류 및 평가}: NIST AI RMF의 4대 기능(거버넌스, 매핑, 측정, 관리)을 기반으로 리스크를 식별하고, 발생 가능성과 영향도 매트릭스를 통해 정량적으로 평가하며, 4T 모델로 대응 전략을 수립한다.

 \item \textbf{정책 준수 모니터링}: KPI 기반 추적, 위반 탐지 및 에스컬레이션, 정기/임시 감사, 대시보드를 통해 정책의 실효성을 보장한다.

 \item \textbf{Policy-as-Code}: OPA 기반 자동화와 CI/CD 파이프라인 통합을 통해 배포 게이트에서 성능, 공정성, 프라이버시, 문서 기준을 자동 검증한다.

 \item \textbf{인시던트 관리}: 4단계 심각도 분류, NIST 기반 5단계 대응 프로세스, 위기 커뮤니케이션 계획을 통해 AI 인시던트에 체계적으로 대응하고, 사후 검토를 통해 조직 학습을 실현한다.
\end{enumerate}

2026년 이후의 환경에서는 정책 문서를 넘어 기술적 통제와 결합된 ``동적 거버넌스''가 필수적이다. Policy-as-Code, 자동화된 영향평가, 실시간 모니터링을 통해 정책이 문서 수준에서 운영 수준으로 전환될 때 비로소 AI 거버넌스는 실질적인 보호 효과를 발휘할 수 있다.
